\documentclass[12pt,a4paper]{article}

\usepackage[italian]{babel}
\usepackage{geometry}
\geometry{margin=2.5cm}
\usepackage{array}
\usepackage{tabularx}
\usepackage{xcolor}

\definecolor{negativeRed}{RGB}{180, 40, 40}

\newcommand{\cognitivetable}[4]{%
	\begin{center}
		\renewcommand{\arraystretch}{1.5}
		\begin{tabularx}{\linewidth}{|X|X|}
			\hline 
			\textbf{Domanda} & \textbf{Osservazioni}
			\\ \hline
			\textbf{1. L'utente sa cosa deve ottenere con l'azione?} (sa cosa deve fare?) & #1 \\ \hline
			\textbf{2. L'utente può individuare lo strumento di interazione?} & #2 \\ \hline
			\textbf{3. L'utente può anche capire che quello è lo strumento giusto per fare ciò che vuol fare?} & #3 \\ \hline
			\textbf{4. Dopo che l'azione è eseguita, l'utente capisce la risposta che ottiene?} & #4 \\ \hline
		\end{tabularx}
	\end{center}
}

\begin{document}
	
	\section*{Documento di Analisi}
	
	\noindent
	\textbf{Documento di valutazione}\\[0.5em]
	\textbf{Codice di valutazione:} 1 \\
	\textbf{Valutatore:} Lublanis \\
	\textbf{Data di valutazione:} 03/08/2026 \\
	\textbf{Sistema da valutare:} Sito web Home Banking \\
	\textbf{Scopo:} Analizzare le difficoltà dell'utente \\
	\textbf{Compito:} \\
	\textbf{Stato iniziale:} \\
	\textbf{Azioni:} 
	
	\vspace{1em}
	
	\subsection*{Valutazione Azione}
	
	% Esempio di utilizzo della macro riutilizzabile:
	\cognitivetable{%
		\textcolor{negativeRed}{Sì e no. L'utente deve capire che per aggiungere una domanda deve modificare il quiz (anche se è ancora vuoto).}%
	}{%
		Sì. Il bottone "modifica quiz" è ben visibile.%
	}{%
		\textcolor{negativeRed}{No. L'utente può non capire che deve selezionare "modifica quiz" essendo il quiz ancora vuoto.}%
	}{%
		Sì. Il sistema mostra la pagina di aggiunta delle domande.%
	}
	
\end{document}